\raggedright
\section{Zero Trust Architecture Principles}

TAC Industries is designed using a Zero Trust security model, where no user, device, or
network segment is implicitly trusted \cite{nist800207}. Instead of relying on an “internal = trusted” perimeter,
access is granted only after explicit verification and policy evaluation. In practice, this drives
three core design decisions in this project: 

(1) Strict network segmentation to reduce lateral movement

(2) Least-privilege access enforced by role and context

(3) Comprehensive logging so actions are attributable and reviewable

A key principle applied throughout the architecture is that reachability does not imply authorization while being connected to a network segment does not grant permission to access sensitive services or data.

\section{Identity as the Security Perimeter}

Modern access decisions are best enforced at the identity layer rather than by network
location alone \cite{nist800207}. TAC Industries treats identity as the primary security boundary: users
authenticate against a central directory, and applications rely on trusted identity assertions
(claims) rather than implementing standalone authentication. This informs the selection of
centralised directory services and federated authentication for web access, and it supports
consistent policy enforcement across both infrastructure and application layers. Role and
attribute information (e.g., group membership and bureau/unit context) can then be used to
restrict what actions are permitted and what data is visible, regardless of how a user connects
to the system.

\section{Privileged Access Management Concepts}

Privileged accounts represent a high-impact risk because compromise can enable full control
of infrastructure \cite{nist800207}. TAC Industries therefore adopts Privileged Access Management concepts
that reduce standing privilege and prevent direct administrator-to-server access from user
endpoints. Administrative activity is routed through a controlled bastion pathway, ensuring
sessions are mediated and logged. This reduces exposure to credential theft and lateral
movement while improving accountability, because privileged actions are tied back to a
named identity and auditable access path rather than unmanaged remote administration.
